\documentclass[11pt]{article}
\usepackage{amsmath,amssymb,amsthm}
\usepackage{bm}
\usepackage[top=1in,bottom=1in,left=1in,right=1in]{geometry}

\newcommand{\E}{\mathbb{E}}
\newcommand{\R}{\mathbb{R}}
\newcommand{\P}{\mathbb{P}}
\newcommand{\Var}{\mathrm{Var}}
\newcommand{\Cov}{\mathrm{Cov}}
\newcommand{\calL}{\mathcal{L}}
\newcommand{\btheta}{\bm{\theta}}
\newcommand{\bxi}{\bm{\xi}}
\newcommand{\bzero}{\bm{0}}
\DeclareMathOperator*{\argmin}{arg\,min}

\newtheorem{assumption}{Assumption}
\newtheorem{proposition}{Proposition}
\newtheorem{remark}{Remark}

\title{Leap complexity: scaling of epochs to escape plateau with batch size}
\author{}
\date{}

\begin{document}
\maketitle

We derive the scaling of the number of epochs to escape a loss plateau with the batch size $B$, under the model: SGD as martingale, diffusion on the plateau, and first passage to a ``hole'' (region of lower loss). The leap complexity is the expected time for this first passage.

\section{Model}

\subsection{SGD as martingale}

Let $\btheta_t \in \R^d$ denote the parameters at step $t$, $\eta$ the learning rate, and $\calL(\btheta) = \frac{1}{n}\sum_{i=1}^n \ell(\btheta; z_i)$ the population loss. With minibatch $\mathcal{B}_t \subset [n]$ of size $|\mathcal{B}_t| = B$, the SGD update is
\begin{equation}
\btheta_{t+1} = \btheta_t - \eta \, \widehat{g}_t, \qquad \widehat{g}_t = \frac{1}{B} \sum_{i \in \mathcal{B}_t} \nabla \ell(\btheta_t; z_i).
\end{equation}
Conditional on $\btheta_t$, the minibatch gradient is unbiased and has variance scaling with $1/B$:
\begin{equation}
\E[\widehat{g}_t \mid \btheta_t] = \nabla \calL(\btheta_t), \qquad \Var(\widehat{g}_t \mid \btheta_t) = \frac{\Sigma(\btheta_t)}{B}, \quad \Sigma(\btheta) = O(1).
\end{equation}
So the noise $\bxi_t = \widehat{g}_t - \nabla \calL(\btheta_t)$ is a martingale difference: $\E[\bxi_t \mid \btheta_t] = \bzero$.

\subsection{Plateau and hole}

On a \emph{plateau}, the full-batch gradient is negligible: $\nabla \calL(\btheta) \approx \bzero$ for $\btheta$ in a flat region $\mathcal{P}$. A ``hole'' is a subset $\mathcal{H} \subset \R^d$ (e.g.\ boundary of $\mathcal{P}$ or a lower-loss basin) such that escaping the plateau means hitting $\mathcal{H}$. We define the \emph{leap} as the event that the trajectory leaves $\mathcal{P}$ and reaches $\mathcal{H}$.

\subsection{Diffusion approximation on the plateau}

On the plateau, the drift $\eta \nabla \calL(\btheta_t)$ is negligible, so one step is
\begin{equation}
\btheta_{t+1} - \btheta_t = -\eta \, \bxi_t, \qquad \E[\bxi_t \mid \btheta_t] = \bzero, \quad \E[\|\bxi_t\|^2 \mid \btheta_t] \asymp \frac{\sigma^2}{B}.
\end{equation}
Over many steps, the rescaled process converges to a diffusion: the displacement over $k$ steps has mean $\bzero$ and variance $\approx k \cdot \eta^2 \sigma^2 / B$. So the effective \emph{diffusion coefficient per step} (in a scalar projection) is
\begin{equation}
D_{\mathrm{step}} = \frac{\eta^2 \sigma^2}{B}.
\end{equation}
Equivalently, the variance of the displacement after one step is $2 D_{\mathrm{step}}$ (in the appropriate direction).

\section{First-passage time (in steps)}

Suppose the trajectory starts at distance $\sim L$ from the hole (in the relevant direction). For a one-dimensional Brownian motion with diffusion coefficient $D$, the mean first-passage time to travel a distance $L$ is $T = L^2 / (2D)$. In our setting, $D = D_{\mathrm{step}} \propto 1/B$, so the mean number of \emph{steps} to escape is
\begin{equation}
T_{\mathrm{steps}} \asymp \frac{L^2}{2 D_{\mathrm{step}}} \asymp \frac{L^2 B}{\eta^2 \sigma^2} \propto B.
\end{equation}
Hence
\begin{equation}
\boxed{\; T_{\mathrm{steps}} \propto B. \;}
\end{equation}
Larger batch size implies smaller noise per step, so more steps are needed to accumulate a displacement of order $L$.

\section{Epochs to escape}

One \emph{epoch} corresponds to $n/B$ steps (one pass over the dataset). So the number of epochs to escape is
\begin{equation}
T_{\mathrm{epochs}} = \frac{T_{\mathrm{steps}}}{n/B} = \frac{T_{\mathrm{steps}} \cdot B}{n}.
\end{equation}
With $T_{\mathrm{steps}} \propto B$, we obtain
\begin{equation}
\boxed{\; T_{\mathrm{epochs}} \propto \frac{B^2}{n}. \;}
\end{equation}
So in this baseline diffusion model, \emph{epochs to escape scale like $B^2$} (with fixed $n$).

\section{Refined scaling: exponent $3/4$}

Empirically, one often observes $T_{\mathrm{epochs}} \propto B^\alpha$ with $\alpha \approx 3/4$ rather than $2$. This can be explained by a bias--variance compromise and a refined first-passage picture.

\subsection{Attempt-rate heuristic}

Instead of a single diffusion coefficient, suppose that in each \emph{epoch} the process has a probability $p(B)$ of ``attempting'' an escape, and that the success of the attempt depends on the noise accumulated during that epoch. The variance of the displacement in one epoch is
\begin{equation}
\sigma^2_{\mathrm{epoch}} = \frac{n}{B} \cdot \frac{\eta^2 \sigma^2}{B} = \frac{n \eta^2 \sigma^2}{B^2}.
\end{equation}
So the typical displacement in one epoch is of order $\sigma_{\mathrm{epoch}} \propto 1/B$. For a fixed barrier at distance $L$, the probability that a single epoch crosses the barrier might scale like a tail probability: e.g.\ if the displacement in one epoch is Gaussian, $\P(\text{escape in one epoch}) \sim \exp(- c B^2)$ for some $c$, which is too small. So we need a different mechanism.

\subsection{Narrow escape / effective dimension}

In a \emph{narrow escape} regime, the hole is a small target (e.g.\ a spherical cap of solid angle $\epsilon$ on the boundary of a ball). The mean first-passage time for a diffusion in $d$ dimensions from the centre of a ball of radius $R$ to a small window of size $\epsilon$ scales like
\begin{equation}
T \sim \frac{R^2}{D} \cdot \frac{1}{\epsilon^{d-1}} \quad \text{(for $d \geq 2$)}.
\end{equation}
If the effective ``size'' of the hole seen by the process depends on $B$ (e.g.\ $\epsilon \propto B^{-\beta}$ because smaller batches have larger fluctuations and ``see'' a larger effective target), then $T \propto B \cdot B^{\beta(d-1)} = B^{1 + \beta(d-1)}$. In \emph{steps}, $T_{\mathrm{steps}} \propto B \cdot B^{\beta(d-1)}$. In \emph{epochs}, $T_{\mathrm{epochs}} = T_{\mathrm{steps}} \cdot B / n \propto B^{2 + \beta(d-1)}$. To get $T_{\mathrm{epochs}} \propto B^{3/4}$, we would need $2 + \beta(d-1) = 3/4$, i.e.\ $\beta(d-1) = -5/4$. For $d=2$, $\beta = -5/4$ (target grows with $B$); for $d=3$, $\beta = -5/8$, etc. So a narrow-escape geometry with an effective target size that grows with batch size can reduce the exponent.

\subsection{Alternative: steps-to-escape scaling}

Suppose instead that the first-passage time in \emph{steps} scales as $T_{\mathrm{steps}} \propto B^\gamma$ for some $\gamma < 1$ (e.g.\ due to correlation of steps or a non-Gaussian tail). Then
\begin{equation}
T_{\mathrm{epochs}} = \frac{T_{\mathrm{steps}} \cdot B}{n} \propto B^{\gamma + 1}.
\end{equation}
For $T_{\mathrm{epochs}} \propto B^{3/4}$ we need $\gamma + 1 = 3/4$, so $\gamma = -1/4$. That would mean $T_{\mathrm{steps}} \propto B^{-1/4}$, i.e.\ \emph{more} steps to escape when $B$ is larger, which is the right direction, but $T_{\mathrm{steps}}$ decreasing with $B$ would require the effective diffusion to \emph{increase} with $B$ (e.g.\ if the number of ``independent attempts'' per step scales with $B$). A more natural interpretation: the \emph{number of steps per epoch} is $n/B$; if the escape is dominated by ``lucky'' epochs in which one or a few steps cause a large displacement, then the relevant quantity might be the number of such steps, which scales like $T_{\mathrm{epochs}} \times (n/B)$. If we posit that the rate of lucky steps per epoch is $\propto 1/B^{1/4}$ (so smaller batches have more lucky steps per epoch), then $T_{\mathrm{epochs}} \propto 1 / (1/B^{1/4}) = B^{1/4}$, which is too small. To get $B^{3/4}$, we could have the rate of lucky steps per epoch $\propto 1/B^{3/4}$, so $T_{\mathrm{epochs}} \propto B^{3/4}$. That is consistent with a bias--variance view: variance per step $\propto 1/B$, and the probability of a ``large'' fluctuation in one epoch might scale like a power of the variance accumulated over the epoch, e.g.\ $(n/B) \cdot (1/B) = n/B^2$ variance per epoch, so a ``typical'' large fluctuation might scale like $1/B$, and the probability of exceeding a threshold could scale like a power of $1/B$, yielding an exponent between 0 and 1.

\subsection{Summary}

\begin{proposition}[Leap complexity scaling]
Under the model (SGD as martingale, diffusion on the plateau with $D_{\mathrm{step}} \propto 1/B$, first passage to the hole):
\begin{itemize}
\item Baseline: $T_{\mathrm{steps}} \propto B$, $T_{\mathrm{epochs}} \propto B^2/n$.
\item Refined (narrow escape or bias--variance): $T_{\mathrm{epochs}} \propto B^\alpha$ with $\alpha \in (0, 2)$. The empirically observed $\alpha \approx 3/4$ is consistent with an effective target size or escape probability that scales with $B$ so that the exponent reduces from $2$ to $3/4$.
\end{itemize}
\end{proposition}

\begin{remark}
The exponent $3/4$ is not derived from first principles here; it is the value observed in experiments (e.g.\ plateau escape for a low-rank network). A full derivation would require specifying the geometry of the plateau and the hole, and solving (or bounding) the first-passage time in that geometry with $D \propto 1/B$.
\end{remark}

\end{document}
